% Chaîne de traitement temps réel et contribution des étapes (figure 4.1)
\begin{tikzpicture}[x=1cm, y=1cm]

  \tikzset{et/.style={fbox, text width=21mm, align=center, minimum height=13mm,
                      font=\sffamily\tiny}}

  % ── Les étapes de la chaîne ───────────────────────────────
  \node[et] (e1) at ( 0.0,0) {\textbf{Capture}\\et transport\\\textit{temps réel}};
  \node[et] (e2) at ( 2.5,0) {\textbf{Détection}\\de fin d'énoncé};
  \node[et] (e3) at ( 5.0,0) {\textbf{Transcription}\\\textit{en continu}};
  \node[et] (e4) at ( 7.5,0) {\textbf{Assemblage}\\du contexte};
  \node[et, fill=figGrisTresClair, densely dashed] (e5) at (10.0,0)
       {\textbf{Récupération}\\documentaire\\\textit{conditionnelle}};
  \node[et] (e6) at (12.5,0) {\textbf{Raisonnement}};
  \node[et] (e7) at (15.0,0) {\textbf{Synthèse}\\vocale};
  \node[et] (e8) at (17.5,0) {\textbf{Restitution}\\au client};

  \foreach \a/\b in {e1/e2, e2/e3, e3/e4, e4/e5, e5/e6, e6/e7, e7/e8}{
    \draw[fluxep] (\a) -- (\b);
  }
  % Contournement lorsque la question n'est pas documentaire
  \draw[flux] (e4.north) -- ++(0,0.85) -- (12.5,0.85)
        node[etiq, above, pos=0.55] {si la demande n'appelle pas de recherche} -- (e6.north);

  % ── Contribution temporelle, en barres ────────────────────
  \def\ybar{-1.35}
  \foreach \x/\w/\lbl in {
      0.0/0.30/{transport},
      2.5/0.55/{250 ms},
      5.0/0.50/{150 à 250},
      7.5/0.20/{30 à 60},
      10.0/0.72/{200 à 350},
      12.5/1.30/{400 à 700},
      15.0/0.50/{150 à 250},
      17.5/0.25/{gigue}}
    {
      \fill[figGrisMoyen, draw=black] (\x-0.42,\ybar) rectangle (\x+0.42,\ybar-\w);
      \node[font=\sffamily\tiny, anchor=north] at (\x,\ybar-\w-0.06) {\lbl};
    }

  \node[font=\sffamily\tiny\itshape, anchor=east, align=right] at (-0.65,\ybar-0.55)
       {Contribution\\au délai perçu};

  \node[note, text width=92mm, anchor=north, align=center] at (8.75,-3.55)
       {Le délai perçu par le client n'est produit par aucune étape en particulier : il est la somme des\\contributions de toutes. Une étape qui dépasse son budget le retire à toutes les autres.};

\end{tikzpicture}
